% !TeX program = xelatex
\documentclass[a4paper,11pt]{ctexart}
\usepackage[margin=25mm]{geometry}
\usepackage{booktabs,tabularx,array,longtable,enumitem,hyperref}
\hypersetup{hidelinks}
\setlength{\parindent}{2em}
\renewcommand{\arraystretch}{1.35}
\title{AI 工具使用详情（核对稿）}
\author{}
\date{2026年8月3日}
\begin{document}
\maketitle

\section*{一、工具名称和版本}
\begin{tabularx}{\textwidth}{>{\bfseries}p{0.25\textwidth}X}
工具名称 & ChatGPT \\
版本／型号 & GPT-5.6 Thinking \\
开发机构 & OpenAI \\
使用日期 & 2026年8月3日（以及参赛队实际使用的其他日期，提交前补全） \\
\end{tabularx}

\section*{二、具体使用目的和环节}
根据当前可见会话，AI 工具至少用于：论文内容转写为 LaTeX、套用数学建模模板、排版优化、
XeLaTeX 编译排错、依据 2026 年官方规范调整电子版格式和检查提交要求。

若 AI 还参与了研究思路、模型建立、参数设定、程序编写、数据处理、数值计算、图表生成或结论分析，
必须在本文件中逐项补充，不能只写“用于润色和排版”。

\section*{三、当前可确认的关键交互记录}
\begin{longtable}{>{\centering\arraybackslash}p{0.07\textwidth}p{0.34\textwidth}p{0.48\textwidth}}
\toprule
序号 & 重要提示词／请求 & AI 回复及采纳、修改情况 \\
\midrule
\endfirsthead
\toprule
序号 & 重要提示词／请求 & AI 回复及采纳、修改情况 \\
\midrule
\endhead
1 & “把第一篇论文按照第二个文件的模板输出对应的 tex。” & 将原论文内容转换为基于 \texttt{cumcmthesis} 的 LaTeX 工程；参赛队应核对正文、公式、图表和参考文献是否与原稿一致。 \\
2 & “按照你的想法，把排版尽可能美观。” & 调整标题、页眉页脚、表格和结论框等视觉样式；后续因竞赛官方格式要求，已撤销可能造成非标准化的自定义页眉、彩色标题和装饰性版式。 \\
3 & “编译有问题。” & 排查主文件、编译器和工程目录结构，统一为 XeLaTeX 编译，并生成可在 Overleaf 使用的完整工程。 \\
4 & “按照数学建模官方的要求。” & 查询全国组委会 2026 年论文格式规范与 AI 工具使用规定，调整为 A4、25 mm 页边距、无目录、摘要页为电子版第一页、页脚居中编号，并增加 AI 使用说明与支撑材料核对项。 \\
\bottomrule
\end{longtable}

\section*{四、采纳和人工修改情况}
当前版本已经采纳的内容包括模板转换、官方版式调整与编译修复。正式提交前，参赛队必须：
\begin{enumerate}[leftmargin=2.5em]
\item 逐项复核模型、参数、公式、程序和数值结果；
\item 补充所有与建模、代码和计算相关的 AI 交互记录；
\item 说明哪些 AI 建议被采纳、部分采纳或未采纳，以及人工修改的具体内容；
\item 确保本文件与论文正文中的 AI 标注、参考文献和实际支撑材料完全一致。
\end{enumerate}

\section*{五、完整性声明}
本文件仅依据当前可见的排版与格式调整会话生成，不能证明已经覆盖此前所有 AI 使用记录。
正式提交前应由参赛队根据完整实际记录补充并重新导出为 \texttt{AI工具使用详情.pdf}。

\end{document}
